\chapter{Typed-Interface Multi-Agent Reinforcement Learning: TI-MARL}
\label{chap:ti-marl}
\markboth{\gls{timarl}}{\gls{timarl}}
\edef\timarlpostdisplaypenalty{\the\postdisplaypenalty}
\postdisplaypenalty=10000

Chapter~\ref{chap:control-benchmark} examined how control and learning methods contribute to the practical use of \gls{marl} in \glspl{rec}. The remaining problem concerns their combined performance: sustaining useful coordination with manageable computation and integration effort, respecting physical and service constraints, and retaining service under disturbance. This chapter develops \gls{timarl} as an integrated response to the scalability, deployability, safety, and robustness requirements defined in Chapter~\ref{chap:introduction}.

The explanation follows a household controller into operation and through changes affecting its resources and measurements. Another charger is installed, a required measurement becomes unavailable, and another building joins the \gls{rec}. These situations connect representation and learning to constraint handling, fault response, and deployment. The chapter develops each mechanism before examining their relationship to the algorithmic approaches reviewed in Chapter~\ref{chap:related-work} and the common comparison in Chapter~\ref{chap:control-benchmark}.

\section{Problem Formulation and Method Overview}
\label{sec:timarl-positioning}
\label{sec:timarl-closest-work}

A useful controller must preserve the relationship between the policy evaluated in development and the decisions applied in an installation. Its inputs must describe the intended resources; its commands must respect their electrical and service conditions; and the loss of information or equipment must leave an explicit path for continued operation or fallback. These responsibilities exist in the original installation and persist as its composition and operating conditions change.

\gls{timarl} connects them through a common typed description used by the learner, local constraint handling, health-dependent action availability, compatibility checks, and execution records. Shared processing and population-aware value estimation support reuse across compatible resources. The method's scientific question is whether these mechanisms together improve continued control across the four research dimensions. An interface whose size can change is one representational capability within that method; the assessment concerns the complete control process.

The neighbouring methods supply substantive algorithmic references. Task-open and compositional learning address coordination and reuse; graph-based energy policies relate resource structure to decisions; safety layers and shielding constrain actions; fault-tolerant training prepares policies for adverse conditions \parencite{anil2025mohito,luketina2025compositional,orfanoudakis2026gnndt,dalal2018safe,mguni2025faulttolerant}. Their mechanisms and supporting processes enter the comparison through the same energy, service, computation, and operational requirements.

\subsection{Typed-Interface Decision Processes}
\label{sec:timarl-decision-processes}

The one-agent reduction isolates interface reconfiguration from cross-agent learning. A \gls{timdp} is defined by

\begin{equation}
\mathcal{M}^{\mathrm{TI}} =
\left\langle
\mathcal{S},\Sigma,\mathcal{H},\mathcal{C},\Gamma^{\mathrm{loc}},P,P_{\Sigma},R,\gamma
\right\rangle,
\label{eq:timdp-tuple}
\end{equation}

where $\mathcal{S}$ is the physical and operational state space; $\Sigma$ and $\mathcal{H}$ are the families of capability schemas and health states; $\mathcal{C}$ is the \gls{tic}; and $\Gamma^{\mathrm{loc}}$ describes local resource relations and constraints. The kernel $P$ governs physical transitions, while $P_{\Sigma}$ governs changes to the interface state, including topology and health. The reward function is $R$, and $\gamma\in[0,1)$ discounts future rewards. Their time-indexed instances, such as $\Sigma_t$ and $\mathcal{H}_t$, describe the current interface. Learning in this one-agent model is denoted \gls{tirl}.

A \gls{timg} adds an active agent population, local information boundaries, coupled outcomes, and shared resources:

\begin{equation}
\mathcal{G}^{\mathrm{TI}} =
\left\langle
\mathcal{S},\mathcal{N},\Sigma,\mathcal{H},\mathcal{C},\Gamma,P,P_{\Sigma},\{R_i\}_{i\in\mathcal{N}},\gamma
\right\rangle.
\label{eq:timg-tuple}
\end{equation}

Here $\mathcal{N}$ is the population of possible agent identities, $\Gamma$ includes local and shared relations, and $R_i$ is the reward function of agent $i$; the remaining objects retain their preceding roles. At time $t$, $\mathcal{N}_t\subseteq\mathcal{N}$ denotes the active population. Each agent $i$ has a schema $\Sigma_{i,t}$ and health state $\mathcal{H}_{i,t}$. The physical state is augmented with the active population and interface state:

\begin{equation}
\tilde{s}_t =
\left(
 s_t,
 \mathcal{N}_t,
 \{\Sigma_{i,t},\mathcal{H}_{i,t}\}_{i\in\mathcal{N}_t},
 \Gamma_t
\right).
\label{eq:timg-augmented-state}
\end{equation}

Here $s_t$ is the physical state, and $\tilde{s}_t$ retains the active members, their schemas and health states, and the current relations $\Gamma_t$. This is an augmented modelling state; each actor observes only its authorised view. The compiler produces the policy-facing interface of agent $i$:

\begin{equation}
\mathcal{I}_{i,t}
=
\mathcal{C}
\left(
\Sigma_{i,t},
\mathcal{H}_{i,t},
\operatorname{view}_i(\tilde{s}_t)
\right)
=
\left(
\mathcal{V}_{i,t},
\mathcal{E}_{i,t},
\mathcal{G}_{i,t}^{a},
\mathcal{P}_{i,t},
\mathcal{B}_{i,t}
\right),
\label{eq:interface-compiler}
\end{equation}

where $\mathcal{V}_{i,t}$ and $\mathcal{E}_{i,t}$ are the typed interface nodes and relations, $\mathcal{G}_{i,t}^{a}$ is the variable set of active action-group instances, $\mathcal{P}_{i,t}$ is the set of valid action-port instances, and $\mathcal{B}_{i,t}$ contains the current parameter domains and local hard constraints. The view function enforces the information boundary: local actors receive their local interface and explicitly authorised aggregate context, not the private observations of other agents.

The policy produces one decision for each active action-group instance:

\begin{equation}
a_{i,t}
=
\left\{
 a_{i,g,t}
\right\}_{g\in\mathcal{G}_{i,t}^{a}},
\qquad
 a_{i,g,t}
=
\left(
 p_{i,g,t},\eta_{i,g,t}
\right),
\label{eq:local-action-bundle}
\end{equation}

where $p_{i,g,t}$ is a selected valid port and $\eta_{i,g,t}$ is its parameter vector. The local action is a variable structured bundle whose contents follow the active interface. A constraint-aware construction $\Phi_i$ maps raw decoded group decisions into a locally feasible bundle:

\begin{equation}
\bar{a}_{i,t}
=
\Phi_i
\left(
 a_{i,t},
 \Gamma^{\mathrm{loc}}_{i,t}
\right),
\qquad
\bar{a}_{i,t}\in\mathcal{A}^{\mathrm{loc}}_{i,t}.
\label{eq:local-feasible-bundle}
\end{equation}

Here $\Phi_i$ is the local construction rule, $\bar a_{i,t}$ is its output, and $\mathcal A^{\mathrm{loc}}_{i,t}$ is the feasible set defined by the current local constraints. The multi-agent joint action is

\begin{equation}
\bar{\mathbf{a}}_t
=
\left(
\bar{a}_{i,t}
\right)_{i\in\mathcal{N}_t}.
\label{eq:joint-action}
\end{equation}

The bold symbol $\bar{\mathbf a}_t$ collects the constructed bundles of the active agents. Each bundle is encoded into resource commands, with execution feedback retaining the action actually applied. The physical transition is

\begin{equation}
s_{t+1}
\sim
P
\left(
\cdot\mid s_t,\bar{\mathbf{a}}_t
\right).
\label{eq:timg-transition}
\end{equation}

The transition kernel $P$ gives the distribution of the next physical state $s_{t+1}$ conditional on $s_t$ and the joint constructed action. Local and \gls{rec} outcomes are coupled through this transition. A representative reward decomposition is

\begin{equation}
r_{i,t}
=
r^{\mathrm{loc}}_{i,t}
+
\alpha_i r^{\mathrm{com}}_t
-
\beta_i c^{\mathrm{loc}}_{i,t}
-
\eta_i c^{\mathrm{shared}}_t,
\label{eq:timarl-reward}
\end{equation}

Here $r_{i,t}$ is the reward of agent $i$, $r^{\mathrm{loc}}_{i,t}$ measures its local outcome, and $r^{\mathrm{com}}_t$ measures the \gls{rec} outcome. The non-negative costs $c^{\mathrm{loc}}_{i,t}$ and $c^{\mathrm{shared}}_t$ represent declared local and shared constraint penalties. The coefficients $\alpha_i$, $\beta_i$, and $\eta_i$ set their relative weights. Local terms can include energy cost and service; \gls{rec} terms can include peak, ramping, renewable utilisation, or a demand-response objective. When $|\mathcal{N}_t|=1$ and $\Gamma_t$ contains only local relations, the model reduces to \gls{tirl}.

The energy formulation uses one building controller per agent. Chargers, batteries, meters, and service sessions are entities within that boundary. Admission concerns new instances or compositions of recognised semantic types. Introducing an unfamiliar physical capability requires its type, interpretation, and control support to be defined before it can participate.

The reduction also clarifies the experimental sequence. \gls{tippo} tests whether the typed local interface and its reconfiguration can be learned without central training. \gls{timappo} keeps the same actor and execution path but uses the active population during value estimation. Differences between the two can therefore be attributed to centralised training rather than to a different deployed controller.

\subsection{Method Overview}
\label{sec:timarl-architecture-overview}



\Cref{fig:timarl-architecture} separates the runtime path from the training path. Each actor evaluation receives one compiled local interface and authorised aggregate signals. It does not receive the raw private state of another member. \gls{timappo}'s critic pools the active agent set during training; \gls{tippo} uses only the local interface. The actor, compiler, projector, and command codec are otherwise the same.

\begin{figure}[H]
\centering
% Editable vector source. Colours distinguish interface, learned policy,
% protection, execution, and the training-only path.
\resizebox{\textwidth}{!}{%
\begin{tikzpicture}[x=1cm,y=1cm,
  stage/.style={draw=#1!65!black,fill=#1!8,rounded corners=3pt,
    text width=2.65cm,minimum height=1.60cm,align=center,font=\small,inner sep=5pt},
  flow/.style={-Latex,line width=0.9pt,draw=black!70},
  train/.style={-Latex,dashed,line width=0.9pt,draw=teal!75!black}]
\draw[rounded corners=5pt,draw=blue!35,fill=blue!2] (-1.65,0.7) rectangle (11.6,-2.35);
% House icon marks a physical controller boundary.
\draw[blue!70!black,line width=1pt] (-1.30,0.13) -- (-.96,.44) -- (-.62,.13);
\draw[blue!70!black,line width=1pt] (-1.23,.16) -- (-1.23,-.2) -- (-.69,-.2) -- (-.69,.16);
\node[anchor=west,font=\small\bfseries,text=blue!70!black] at (-.4,.15) {Building A: local decision path};
\node[anchor=east,font=\scriptsize,text=black!70] at (11.3,.15) {same actor weights in compatible buildings};
\node[stage=cyan] (compile) at (0,-1.2) {\textbf{1\quad Compile}\\Entity bindings\\Live interface};
\node[stage=orange] (actor) at (3.3,-1.2) {\textbf{2\quad Decide}\\Shared actor\\Mode / fraction};
\node[stage=violet] (project) at (6.6,-1.2) {\textbf{3\quad Construct}\\Local projector\\Feasible bundle};
\node[stage=gray] (codec) at (9.9,-1.2) {\textbf{4\quad Encode}\\Asset commands\\Execution record};
\draw[flow] (compile) -- (actor);
\draw[flow] (actor) -- (project);
\draw[flow] (project) -- (codec);
\node[draw=blue!35,fill=blue!3,rounded corners=3pt,align=center,text width=3.4cm,minimum height=1.05cm,font=\small] (others) at (.25,-3.8) {\textbf{Other active buildings}\\Separate local paths};
\node[draw=gray!55,fill=gray!8,rounded corners=3pt,align=center,text width=5.8cm,minimum height=1.05cm,font=\small] (env) at (8.1,-3.8) {\textbf{Shared energy environment}\\Applied actions, next measurements,\\local rewards and community outcomes};
\draw[flow] (codec.south) -- (9.9,-3.23);
\draw[flow] (others.east) -- node[above,font=\scriptsize] {local commands} (env.west);
\draw[flow] (env.east) -- (11.85,-3.8) -- (11.85,1.05) -- (-2,1.05) -- (-2,-1.2) -- (compile.west);
\node[font=\scriptsize,fill=white,inner sep=2pt] at (6.6,1.05) {new measurements and events};
\node[draw=teal!65!black,fill=teal!7,dashed,rounded corners=4pt,text width=9cm,align=center,minimum height=1.08cm,font=\small] (critic) at (5.0,-5.65) {\textbf{Training only}\\\gls{tippo}: local critic \qquad \gls{timappo}: central set critic};
\draw[train] (env.south) -- (8.1,-5.06);
\draw[train] (others.south) -- (.25,-5.65) -- (critic.west);
\draw[train,preaction={draw=white,line width=3pt}] (3.3,-5.06) -- node[right,pos=.68,font=\scriptsize,fill=white,inner sep=2pt] {actor updates} (actor.south);
\end{tikzpicture}%
}

\caption{\gls{timarl} decision path within each building. Local construction precedes execution; the \gls{rec} critic participates only in training.}
\label{fig:timarl-architecture}
\end{figure}

Table~\ref{tab:timarl-components} identifies the role of each component. The same interface snapshot accompanies the decision through these stages, so a charger retains its identity when its observations become actor inputs and its selected power becomes an equipment command.

\begin{table}[htbp]
\centering
\caption{Main \gls{timarl} components and their interpretation in the running energy example.}
\label{tab:timarl-components}
\begin{tabularx}{\textwidth}{@{}p{0.22\textwidth}p{0.38\textwidth}X@{}}
\toprule
\textbf{Component} & \textbf{Generic role} & \textbf{Energy example} \\
\midrule
\gls{tic} & Binds deployment-neutral member contracts to current entities, observations, events, health, groups, bounds, and constraints & Adds a known charger group when it becomes active; scopes failure to the affected group; contracts \gls{v2g} when its evidence is stale \\
Hierarchical typed encoder & Aggregates observations by channel and sensor before forming a local latent state & Keeps two charger sessions distinct through semantic and structural identities without learning an embedding for a concrete charger identifier \\
Shared hybrid actor & Produces one categorical mode and one bounded continuous fraction for each active group & Controls two chargers, one battery, and one deferrable service with shared group-type heads rather than an asset-specific output head \\
Analytic local projector & Applies validity, dynamic bounds, service floors, and joint total/per-phase headroom & Preserves an urgent \gls{ev} or mandatory deferrable start while scaling competing charging decisions to the available connection capacity \\
Local or central set critic & Supplies \gls{tippo} local values or \gls{timappo} values conditioned on the active member set & Retains the same decentralised actor while comparing independent and \gls{rec}-aware value estimation \\
\bottomrule
\end{tabularx}
\end{table}

The architecture is deliberately asymmetric between training and operation. The actor remains local in both phases. Centralisation is introduced only in \gls{timappo}'s value function and the experiment infrastructure, where it can improve credit assignment without creating a runtime dependency on unrestricted global state. \gls{rec} observations remain typed observation parts with explicit scope, provenance, time, and health.

\section{Typed Interfaces and Runtime Compilation}
\label{sec:runtime-interface-model}

The interface provides a stable vocabulary for a changing installation. A schema declares what a resource means and which capabilities it can expose; compilation binds that declaration to the resources and evidence present at the current step. This separates a new instance of a familiar charger from a new kind of equipment requiring additional control semantics.

\subsection{Schemas, Entities, and Dynamic Module Instances}
\label{sec:timarl-entities}

A capability schema describes the potential interface of an agent. It identifies the agent role, module types, observation templates, action-group templates, constraints, compatibility requirements, and fallback semantics. Runtime state identifies which module instances and controlled entities are active now. Separating the two prevents fast-changing values such as \gls{soc} or connection state from being treated as schema revisions.

A building can contain several instances of the same module type. In the running example, agent $A$ contains two charger instances, one stationary battery, one deferrable-service module, and one main meter:

\begin{equation}
\mathcal{M}_{A,t}
=
\{m^{\mathrm{meter}}_1,
  m^{\mathrm{charger}}_1,
  m^{\mathrm{charger}}_2,
  m^{\mathrm{battery}}_1,
  m^{\mathrm{def}}_1\}.
\label{eq:running-modules}
\end{equation}

Here $\mathcal M_{A,t}$ is the module set of building $A$ at time $t$; the superscripts identify module types and the subscripts distinguish their instances. Each charger can be related to a currently connected \gls{ev} session. The following observation identifies the \gls{soc} source for one such session:

\begin{quote}\small
\begin{minipage}{\linewidth}
\texttt{semantic type: EV\_SOC}\\
\texttt{entity: ev\_session\_17}\\
\texttt{source module: charger\_1}
\end{minipage}
\end{quote}

The source and session identities distinguish this measurement from the \gls{soc} of the vehicle connected to \texttt{charger\_2}. The corresponding action group is similarly bound:

\begin{quote}\small
\begin{minipage}{\linewidth}
\texttt{group type: EV\_CHARGER\_MODE}\\
\texttt{owner module: charger\_1}\\
\texttt{target entity: ev\_session\_17}
\end{minipage}
\end{quote}

The compiler represents these relationships as typed edges such as \texttt{owns}, \texttt{connected\_to}, \texttt{observes}, \texttt{controls}, and \texttt{requires}. Entity binding is encoded in the interface semantics and does not depend on position in a flat vector.

When \texttt{charger\_3} is installed, a module-addition event creates a new module node, its observation parts, one \gls{ev}-charger action-group instance, valid port instances, and the relevant constraint edges. The neural network is not extended with a new permanent head. The same group and port decoder used for the existing chargers is applied to the new instance. Removing the charger removes that subgraph. A semantically unknown module is admitted only when a compatible parent type and decoder are available; otherwise its learned control path remains inactive and fallback semantics apply.

The implementation deliberately separates four identifiers that flat action layouts often conflate. A member owns the local controller boundary; a charger identifies installed infrastructure; an \gls{ev} identifies a vehicle; and a session records the temporary occupancy of one charger by one \gls{ev}. A domestic charger can retain the same vehicle over many sessions, whereas one office charger can serve several vehicles in one day. The action group belongs to the physical charger, while service observations belong to the currently bound session. Session replacement therefore rebinds runtime values without rebuilding an unchanged charger structure.

\subsection{Typed Observation Parts and Community Context}
\label{sec:timarl-observations}

An observation part is a typed object with an identity and declared semantics. A representative definition is

\begin{equation}
o_k
=
\left(
\tau_k,
\mathrm{id}_k,
\mathrm{entity}_k,
\mathrm{source}_k,
 x_k,
 u_k,
 q_k,
 h_k,
 t_k,
 s_k
\right),
\label{eq:observation-part}
\end{equation}

where $k$ indexes an observation part, $\mathrm{id}_k$ identifies that part, $\mathrm{entity}_k$ identifies the entity described, and $\mathrm{source}_k$ identifies its source. The remaining entries are $\tau_k$, the semantic type, $x_k$ the value, $u_k$ the unit, $q_k$ quality or uncertainty metadata, $h_k$ the health state, $t_k$ the timestamp, and $s_k$ the scope. The scope can be self, local, neighbourhood, \gls{rec}, or global. The source can be a physical sensor, derived signal, forecast, external platform, or coordinator message.

\gls{rec} information is represented through the same object model. A building may receive numerical aggregates such as \gls{rec} net demand, transformer loading, remaining headroom, \gls{pv} surplus, a price multiplier, or a demand-response target. These values do not reveal the raw private state of another building. Their provenance and health remain explicit:

\begin{quote}\small
\begin{minipage}{\linewidth}
\texttt{semantic type: COMMUNITY\_TRANSFORMER\_LOADING}\\
\texttt{scope: community; source: platform\_service}\\
\texttt{unit: percent; health: HEALTHY}
\end{minipage}
\end{quote}

In the energy instantiation, an \gls{ev} connected from 08:00 until 18:00 with a current \gls{soc}, target \gls{soc}, maximum charging power, and departure requirement already exposes the information from which a policy can learn temporal flexibility. Useful derived quantities such as time to departure, remaining required energy, or required average charging power can be supplied as typed observation parts when they improve learning or interpretability. Their inclusion is fixed by the experimental configuration.

\Cref{tab:timarl-running-observations} lists representative observation parts for the running example.

\begin{table}[htbp]
\centering
\small
\caption{Representative typed observation parts for a building with two chargers. Values are illustrative.}
\label{tab:timarl-running-observations}
\begin{tabularx}{\textwidth}{@{}p{0.19\textwidth}p{0.16\textwidth}p{0.13\textwidth}p{0.12\textwidth}X@{}}
\toprule
\textbf{Semantic type} & \textbf{Entity/source} & \textbf{Value} & \textbf{Scope} & \textbf{Decision meaning} \\
\midrule
House net power & main meter & 2.1 kW & self & Current aggregate import used to evaluate local electrical limits \\
\gls{ev} \gls{soc} & session 17 / charger 1 & 0.35 & local & Energy state of the \gls{ev} controlled through charger 1 \\
\gls{ev} target \gls{soc} & session 17 & 0.80 & local & Minimum service target at departure \\
Departure time & session 17 & 10 h & local & Temporal slack available to the charging policy \\
\gls{ev} \gls{soc} & session 28 / charger 2 & 0.62 & local & Distinct \gls{soc} bound to charger 2 and its connected session \\
Battery \gls{soc} & battery 1 & 0.70 & local & Available storage state \\
\gls{rec} headroom & platform service & 5.5 kW & \gls{rec} & Aggregate remaining shared capacity made visible to all authorised agents \\
\gls{rec} peak state & platform service & high & \gls{rec} & \gls{rec} context used by local actors and the reward \\
\bottomrule
\end{tabularx}
\end{table}

\subsection{Dynamic Action Groups and Typed Parameterised Ports}
\label{sec:timarl-action-ports}

The output interface is organised around action-group instances. Each controllable module creates one or more groups whose ports are mutually exclusive within the group but can be combined with ports from other groups. For example,

\begin{align}
\mathcal{P}_{\mathrm{charger}_1,t}
&=\{\textsc{idle-ev},\textsc{charge-ev}(p),\textsc{discharge-ev}(p)\},\\
\mathcal{P}_{\mathrm{battery}_1,t}
&=\{\textsc{idle-battery},\textsc{charge-battery}(p),\textsc{discharge-battery}(p)\},\\
\mathcal{P}_{\mathrm{def}_1,t}
&=\{\textsc{idle},\textsc{start}\}.
\end{align}

Each set $\mathcal P$ lists the ports of the named resource at time $t$. In these examples, the argument $p$ denotes a requested physical power; the deferrable service has a discrete start decision. An action port itself is a typed callable capability:

\begin{equation}
p
=
\left(
\tau_p,
\mathrm{id}_p,
\mathrm{group}_p,
\mathrm{owner}_p,
\Theta_p,
U_p,
B_{p,t},
V_{p,t},
D_p,
K_p
\right),
\label{eq:action-port}
\end{equation}

where $\tau_p$ is the port type, $\mathrm{id}_p$ its identity, $\mathrm{group}_p$ the action group to which it belongs, and $\mathrm{owner}_p$ its owning module. The remaining entries are $\Theta_p$, the parameter schema, $U_p$ the units, $B_{p,t}$ the current parameter domain, $V_{p,t}$ the validity state, $D_p$ the observation and capability dependencies, and $K_p$ the resource effects used by local constraints, critics, and optional safety checks.

A new charger creates another instance of the charger action group. It does not create another set of network parameters. The implemented actor emits a categorical mode and a Beta-distributed fraction in $[0,1]$ for every active group. The fraction is relative to the currently compiled physical upper bound and is converted to a signed simulator command only once by the codec. This convention avoids applying a health-dependent or topology-dependent bound twice. A deferrable start is indivisible: after \textsc{start} is selected, its typed cycle profile determines the subsequent demand rather than a continuous power fraction.

\begin{table}[htbp]
\centering
\small
\caption{Representative action ports for two charger instances and one stationary battery.}
\label{tab:timarl-running-ports}
\begin{tabularx}{\textwidth}{@{}p{0.17\textwidth}p{0.17\textwidth}p{0.16\textwidth}p{0.22\textwidth}X@{}}
\toprule
\textbf{Port} & \textbf{Owner/target} & \textbf{Parameter} & \textbf{Current validity and bounds} & \textbf{Dependencies and local effects} \\
\midrule
Charge \gls{ev} & charger 1 / session 17 & power [kW] & valid, $[0,7.2]$ & connected session; increases house import \\
Discharge \gls{ev} & charger 1 / session 17 & power [kW] & valid, $[0,4.0]$ & \gls{v2g}, trustworthy \gls{soc}, reserve; reduces house import \\
Charge \gls{ev} & charger 2 / session 28 & power [kW] & valid, $[0,11.0]$ & distinct connection and \gls{ev} state; increases house import \\
Discharge \gls{ev} & charger 2 / session 28 & power [kW] & absent & charger 2 has no \gls{v2g} capability \\
Discharge battery & battery 1 & power [kW] & valid, $[0,3.0]$ & battery \gls{soc} and reserve; reduces house import \\
\bottomrule
\end{tabularx}
\end{table}

For example, suppose \texttt{charger\_1} is rated at 7.4~kW but the current connection and \gls{ev} state reduce its compiled charging bound to 4.0~kW. A sampled fraction of 0.75 requests 3.0~kW. The network output remains in the same numerical domain when the available bound later changes; the interface gives the output its current physical meaning. The trace stores the selected mode, fraction, compiled bound, requested command, projected command, and applied simulator value separately.

\section{Policy Architecture and Multi-Agent Learning}
\label{sec:tip-and-learning}
\label{sec:timarl-multiagent-learning}

Once the interface has been compiled, the actor must interpret its current observations and produce a decision for each controllable group. Shared processing allows another known charger to reuse the existing policy. Centralised training can then account for the interaction between buildings while preserving the local information available during operation.

\subsection{Local Information and Hierarchical Encoding}
\label{sec:timarl-local-info}

Each agent executes a logically local \gls{tip}. A building actor receives its compiled interface, optional memory, and authorised aggregate \gls{rec} observations. It does not receive the raw loads, \gls{ev} states, or private service requirements of other buildings. This boundary is identical during evaluation and deployment; privileged information used by a central critic during training is not silently passed to the actor.

The current implementation shares one actor across all compatible members:

\begin{equation}
\pi_i
=
\pi_{\theta}
\left(
\mathcal{I}_{i,t},m_{i,t}
\right),
\label{eq:role-sharing}
\end{equation}

Here $\pi_i$ is the policy instance of agent $i$, $\pi_\theta$ is the shared actor with parameters $\theta$, $\mathcal I_{i,t}$ is its local snapshot, and $m_{i,t}$ is optional local memory. Parameter sharing preserves separate observations. House $A$ and House $B$ execute the same weights on different local snapshots, and member role remains an encoded feature. The compiler rejects unknown semantic or action-group types rather than silently allocating a random representation. This keeps zero-shot admission limited to new instances and compositions of known contracts.

This structure supports population changes without a permanent actor slot per agent. A new building instantiates the appropriate local policy family after schema compatibility has been checked. Removing a building removes its actor instance from the active population; the shared parameter set and the actors of the remaining agents are unchanged.

\label{sec:timarl-local-encoder}

The compiled local interface is encoded as variable-cardinality typed sets. The hierarchy follows the simulator's measurement structure: observation parts are grouped into channels, channels into sensors, and sensors into the local member representation. Structural relations remain available when an action group asks for its relevant observations. This hierarchy follows the implemented measurement structure and preserves the binding needed by the decoder.

Each observation token combines its normalised value with quality and validity indicators, health, semantic type, sensor type, channel type, unit, scope, declared use, a deterministic semantic fingerprint, and temporal-horizon features. The first mapping is

\begin{align}
z^{o}_{k}
&=
E^{o}_{\theta}
\left(
x_k,q_k,h_k,e_{\tau_k},e_{\mathrm{sensor}},e_{\mathrm{channel}},
e_{u_k},e_{s_k},e_{\mathrm{use}},f^{\mathrm{id}}_k,f^{\mathrm{time}}_k
\right).
\end{align}

Here $z^o_k$ is the embedding of observation part $k$, and $E^o_\theta$ is its learned encoder. The values $x_k$, $q_k$, and $h_k$ retain the meanings in Equation~\ref{eq:observation-part}; each $e$ is an embedding of the named metadata field. The features $f^{\mathrm{id}}_k$ and $f^{\mathrm{time}}_k$ encode semantic identity and temporal information. Permutation-equivariant relational layers aggregate observation tokens within their channels and sensors, followed by learned pooling over the local sensor set:

\begin{equation}
z_{i,t}
=
F_{\theta}
\left(
\left\{\left\{\{z^o_k\}_{k\in c}\right\}_{c\in s}\right\}_{s\in i}
\right).
\label{eq:tip-encoder}
\end{equation}

The mapping $F_\theta$ aggregates the nested sets: $k$ indexes observation parts, $c$ channels, and $s$ sensors belonging to member $i$. The common latent $z_{i,t}$ summarises the complete local decision context. Concrete asset identifiers are not learned as permanent embeddings. Instance distinction is carried by the compiled structural binding and deterministic observation identity, while semantic families determine reusable parameters. Reordering observations, chargers, or members therefore does not change the intended meaning.

Each active action group starts from its type embedding and the local latent. In the action-conditioned variant, the group also queries observation tokens related to its owner, target entity, and dependencies. The two chargers can use the same mode and fraction heads while attending to different sessions. Attention across the active group set then exposes local competition before decoding. Adding a charger adds one group evaluation; it does not change the actor or critic parameter count.

During replay and \gls{bc} updates, immutable type indices and set membership are prepared once and reused across optimisation epochs. Runtime values and health are never cached as static input. Packed evaluation removes repeated per-agent device transfers while preserving the same set semantics. These changes affect throughput, not the mathematical information surface.

\subsection{Dynamic Grouped Action Decoding}
\label{sec:timarl-action-decoder}

The actor does not contain a permanent neural head for each asset instance. A group-type head is applied to every active group of that type. For group $g$ and candidate mode $m$,

\begin{equation}
\ell_{i,g,m,t}
=
f^{\mathrm{select}}_{\theta}
\left(
 z_{i,t},z_{i,g,t},e_{\tau_g}
\right),
\label{eq:port-logit}
\end{equation}

Here $\ell_{i,g,m,t}$ is the logit for candidate mode $m$ of group $g$, $z_{i,g,t}$ is its contextual embedding, $e_{\tau_g}$ represents its type, and $f^{\mathrm{select}}_\theta$ is the shared selection head. The categorical mode is selected after invalid modes have been masked:

\begin{equation}
\Pr
\left(
 m_{i,g,t}=m
\right)
=
\operatorname{Softmax}_{m\in\mathcal{M}^{\mathrm{valid}}_{i,g,t}}
\left(
\ell_{i,g,m,t}
\right).
\label{eq:port-selection}
\end{equation}

Here $m_{i,g,t}$ is the selected mode, $\mathcal M^{\mathrm{valid}}_{i,g,t}$ is its admissible candidate set, and Softmax normalises the corresponding logits into probabilities. A second head generates positive Beta parameters. Sampling $u_{i,g,t}\sim\operatorname{Beta}(\alpha_{i,g,t},\beta_{i,g,t})$ gives the bounded fraction

\begin{equation}
\eta_{i,g,t}
=
u_{i,g,t},
\qquad
(\alpha_{i,g,t},\beta_{i,g,t})
=
d_{\tau_g}(z_{i,g,t}).
\label{eq:port-parameters}
\end{equation}

The decoder $d_{\tau_g}$ produces positive shape parameters $\alpha_{i,g,t}$ and $\beta_{i,g,t}$ for the Beta distribution; its sample $u_{i,g,t}$ supplies the fraction $\eta_{i,g,t}$. The log probability used by \gls{ppo} is the sum of the categorical and Beta terms. At training time the two parts are sampled. Deterministic replay can use categorical argmax or an \texttt{expected\_signed} rule that combines charge and discharge probabilities, the Beta mean, a configurable gain, a deadband, and a non-idle logit margin. These controls were added because an apparently small stochastic preference can become an implausibly persistent deterministic charge or discharge command after export.

All groups are decoded as one bundle from the same local context. A representative raw bundle is

\begin{equation}
a^{\mathrm{raw}}_{A,t}
=
\left\{
\begin{array}{l}
\texttt{charger\_1}:\textsc{charge-ev}(4.0\ \mathrm{kW}),\\
\texttt{charger\_2}:\textsc{charge-ev}(3.5\ \mathrm{kW}),\\
\texttt{battery\_1}:\textsc{discharge-battery}(1.0\ \mathrm{kW}),\\
\texttt{deferrable\_1}:\textsc{start}
\end{array}
\right\}.
\label{eq:running-raw-bundle}
\end{equation}

Here $a^{\mathrm{raw}}_{A,t}$ collects the four proposed decisions for building $A$ at time $t$; their powers are illustrative. The bundle-level representation allows the actor to learn priorities among the two \gls{ev} sessions, battery use, a pending service, and \gls{rec} conditions. It does not by itself guarantee that the combined power remains within local electrical limits or that an urgent service remains feasible. Those requirements are handled by the final local construction.

\subsection{Centralised Training and Local Execution}
\label{sec:timarl-joint-training}

\gls{timarl} does not require a \gls{rec} controller agent. Each building acts independently from its own compiled interface and authorised \gls{rec} observations. The method becomes multi-agent because the local bundles form one joint action, the physical transition depends on that joint action, local actions compete for shared resources, and rewards include local and \gls{rec} consequences.

At each training step, the constructed local bundles form the joint action in Equation~\ref{eq:joint-action}. Its transition produces the next interfaces, local rewards, and \gls{rec} outcomes. Simultaneous charging can reduce shared headroom even when both buildings satisfy their own connection limits. Repeated interaction supplies the experience from which local actors learn those collective consequences.

Joint training means that actor updates are derived from the same joint transition. It does not mean that actors share private runtime inputs or that one global actor replaces them. Stable member identities govern the rollout when population changes. A joining member has no predecessor in Generalised Advantage Estimation; a departing member receives a terminal boundary; members that remain can bootstrap across the topology event. The rollout is not cleared and the networks are not resized.

\label{sec:timarl-central-critic}

A conventional central critic often concatenates all agent observations. That representation is tied to a fixed order and population. \gls{tippo} avoids centralisation and estimates $V_i$ from the local typed snapshot. \gls{timappo} retains the same actor but estimates values from the variable active set.

The implemented central critic encodes each member with the same typed snapshot encoder, mean-pools the active local latents into a \gls{rec} summary, and evaluates each stable member identity from its local and \gls{rec} representations:

\begin{equation}
V_{\phi,i,t}
=
g_{\phi}
\left(
z_{i,t},
\frac{1}{|\mathcal{N}_t|}
\sum_{j\in\mathcal{N}_t}z_{j,t}
\right).
\label{eq:graph-critic}
\end{equation}

Here $V_{\phi,i,t}$ is the value estimate for member $i$, $g_\phi$ is the learned value mapping, and $z_{j,t}$ is the encoded snapshot of active member $j$. The mean contains all $|\mathcal N_t|$ active members. The mapping is permutation equivariant: reordering members reorders their values but does not alter their association. Cardinality can change without reconstructing a joint input layer. Mean pooling supplies a compact \gls{rec} context. Distinct feeder neighbourhoods with identical aggregate summaries receive the same pooled representation. Representing those neighbourhoods explicitly would require a richer critic input or aggregation rule.

\subsection{Proximal Policy Optimisation with Typed Interfaces}
\label{sec:timarl-backbones}

\gls{timarl} defines the interface, actor structure, local bundle semantics, and population handling. The current learner supports two controlled \gls{ppo} variants.

\paragraph{TI-PPO.}
The shared actor and a parameter-shared local critic are updated from local member trajectories. This is the typed one-agent learner applied independently across the active population. It isolates the effect of the interface and actor from centralised value estimation.

\paragraph{TI-MAPPO.}
The actor is paired with the central set critic in \cref{eq:graph-critic}. The critic may use the active joint training context; the actor retains its local information boundary. A checkpoint pins the backbone and critic kind, so a \gls{tippo} state cannot be loaded into \gls{timappo} by accident.

For a transition generated by old parameters $\theta_{\mathrm{old}}$, the hybrid group probability ratio is

\begin{equation}
r_{i,g,t}(\theta)
=
\exp\left[
\log\pi_{\theta}(m_{i,g,t},\eta_{i,g,t}\mid\mathcal{I}_{i,t})
-
\log\pi_{\theta_{\mathrm{old}}}(m_{i,g,t},\eta_{i,g,t}\mid\mathcal{I}_{i,t})
\right].
\label{eq:timarl-ratio}
\end{equation}

Here $r_{i,g,t}(\theta)$ compares the current and collection-policy probabilities for the same recorded mode, fraction, and interface. The parameters $\theta_{\mathrm{old}}$ identify the collection policy. With joint-agent credit, the clipped \gls{ppo} objective uses the sum of group log probabilities for that member. With typed-group credit, the ratio and advantage remain attached to group $g$:

\begin{equation}
L^{\mathrm{group}}(\theta)
=
\mathbb{E}_{i,g,t}
\left[
\min\left(
r_{i,g,t}(\theta)\hat{A}_{i,g,t},
\operatorname{clip}(r_{i,g,t}(\theta),1-\epsilon,1+\epsilon)
\hat{A}_{i,g,t}
\right)
\right].
\label{eq:timarl-group-ppo}
\end{equation}

Here $\hat A_{i,g,t}$ is the advantage estimate assigned to group $g$, $\epsilon$ is the clipping tolerance, and the expectation averages recorded agent--group transitions. The operator $\operatorname{clip}$ bounds the ratio to $[1-\epsilon,1+\epsilon]$. The group critic supplies a matching baseline. \gls{ev} service and charger-constraint consequences are routed to the responsible charger family; stationary-storage consequences reach the storage group; deferrable completion reaches its start group. Economic settlement and genuinely shared grid effects remain common. This avoids teaching the second charger that it caused a service failure belonging to the first, while preserving \gls{rec} coupling.

The comparison with \gls{matd3} in Chapter~\ref{chap:control-benchmark} retains a practical reference from a different learning family. Attribution to the interface is assessed more directly through \gls{ppo}-based controls with matched actor inputs, training budgets, and protection.

\subsection{Training Data and Time Resolution}
\label{sec:timarl-learning-engineering}

Moving the controller to 15-minute data increased the effective learning horizon and exposed numerical failures that were less visible in the earlier hourly experiments. \gls{timarl} therefore treats the physical duration of a transition as part of the learner contract. When $\gamma_{\mathrm{ref}}$ and $\lambda_{\mathrm{ref}}$ are declared for reference duration $\Delta t_{\mathrm{ref}}$, the factors used at runtime are

\begin{equation}
\gamma_{\Delta t}=\gamma_{\mathrm{ref}}^{\Delta t/\Delta t_{\mathrm{ref}}},
\qquad
\lambda_{\Delta t}=\lambda_{\mathrm{ref}}^{\Delta t/\Delta t_{\mathrm{ref}}}.
\label{eq:timarl-timebase}
\end{equation}

Here $\Delta t$ is the transition duration and $\Delta t_{\mathrm{ref}}$ the reference duration. The factors $\gamma$ and $\lambda$ control reward discounting and Generalised Advantage Estimation respectively, preserving their declared time basis as the step changes. The implementation normalises critic targets, uses a Huber value loss by default, clips policy ratios and gradient norms, and records approximate \gls{kl} divergence, clipping fraction, explained variance, and finite-gradient checks. Advantages can be normalised globally or per stable member identity. Entropy coefficients can differ by action-group family; a charger and a deferrable start do not need the same exploration pressure. A policy anchor can constrain movement from a selected actor and can be reset deliberately when a checkpoint is resumed.

The rollout is on-policy and retains the raw typed decision, projected decision, intervention scope, stable agent and group identities, reward decomposition, termination semantics, current and next snapshot hashes, and simulator execution feedback. Teacher actions and auxiliary-planning examples never enter that \gls{ppo} rollout. This boundary matters: an update remains attributable to experience generated by the current policy even when demonstrations are used elsewhere in training.

\label{sec:timarl-prior-knowledge}

The use of prior knowledge in Section~\ref{sec:control-ablations} motivates an optional \gls{bc} warm start with \gls{rbc} Smart Local (the implementation class \texttt{RBCSmartPolicy}). The teacher acts on the physical simulator interface; its requested commands are decoded back into the active typed groups and stored with the exact interface snapshot on which they were valid. Pretraining updates the actor only. The critic is not trained from teacher returns, and the complete demonstration snapshots are discarded after pretraining.

Idle actions dominate many annual trajectories. A plain average cross-entropy loss can therefore learn an impressive cloning accuracy while missing the charge, discharge, \gls{v2g}, or start events that make the teacher useful. The implementation supports bounded inverse-frequency weighting and a hierarchical mode mean. It can run a final unweighted calibration phase so that balancing rare actions does not permanently distort the deployed prior.

Two optional causal auxiliaries were added after the warm-start experiments exposed sparse service decisions. The \gls{ev} planner labels typed states using only current service variables and declared forecasts: urgent charging, price-aware delay, authorised \gls{v2g} with departure reserve, or idle. The storage planner identifies \gls{pv} capture, low-price charging, high-price discharge, and idle from the same deployment-visible surface. Each auxiliary predicts both mode and fraction, trains the actor only, and keeps a bounded reservoir stratified by decision reason. Rare urgent or arbitrage examples therefore survive without contaminating the \gls{ppo} rollout.

These planners supply training labels; runtime decisions remain with the actor and analytic projector. Pure \gls{timappo}, \gls{bc} warm start, \gls{ev} auxiliary, storage auxiliary, and their combinations remain separate ablations. This separation makes the contribution of prior knowledge testable: prior knowledge is attached to a typed and causal state, while the learned actor still owns execution.

\section{Constraint Handling and Degraded Operation}
\label{sec:tcg-safety}

The compiled interface serves two connected purposes in protection. It determines which decisions remain meaningful under the current evidence, then supplies the physical and service constraints used to construct their joint local effect. A charger fault and a building power limit therefore enter different stages of the same decision path.

\subsection{Health-Aware Causal Compilation}
\label{sec:timarl-health}

Health is part of the interface semantics and is also available to the neural network. The implemented states are \textsc{healthy}, \textsc{degraded}, \textsc{stale}, \textsc{missing}, \textsc{failed}, and \textsc{unknown}. Fault evidence remains a separate record. A simulator event such as \texttt{dropout} or \texttt{stuck} is not mapped to one universal health label; the compiler considers event duration, freshness, criticality, declared availability, channel semantics, and recovery hysteresis.

A charger failure disables only the associated charger group and its ports. Other modules controlled by the member remain active. A stale \gls{ev} \gls{soc} can preserve charging while blocking \gls{v2g} because the discharge reserve can no longer be established. Failure of the main house meter has a wider consequence: the complete member is isolated when aggregate import or phase headroom cannot be verified. Loss of a \gls{rec} signal removes that coordination input but leaves the local controller available. During a grid outage, charging, \gls{v2g}, and a new deferrable start are blocked; safe stationary-storage discharge may remain available when its local dependencies still hold.

The compiler computes the transitive effect of a health event through a dependency closure:

\begin{equation}
\mathcal{H}^{\star}_{i,t}
=
\operatorname{Closure}
\left(
\mathcal{H}_{i,t},D_{i,t}
\right),
\label{eq:health-closure}
\end{equation}

where $D_{i,t}$ is the declared dependency relation, $\operatorname{Closure}$ propagates the observed health state $\mathcal H_{i,t}$ through it, and $\mathcal H^\star_{i,t}$ records the resulting consequences. These consequences concern observations, ports, groups, bounds, local constraints, and operational modes. For the \gls{ev} \gls{soc} example, the closure can be summarised as

\begin{equation}
\mathcal P^{\mathrm{valid}}_{i,t}
\leftarrow
\mathcal P^{\mathrm{valid}}_{i,t}\setminus
\{p:\tau_p=\textsc{discharge-ev},\ \mathrm{owner}_p=\texttt{charger\_1}\}.
\label{eq:timarl-health-port-removal}
\end{equation}
Here $\mathcal P^{\mathrm{valid}}_{i,t}$ is the valid port set of member $i$, $\tau_p$ is a port's type, and $\mathrm{owner}_p$ identifies its equipment. Set subtraction removes only the discharge capability of the affected charger. Charging remains available only if its own dependencies and service constraints can still be established.

The dependency rule changes the decision interface itself, giving the policy an explicit account of the capabilities still available. Recovery is explicit: a group or port returns only after its declared fresh-data and hysteresis conditions hold. The closure log records the subject, state, rule, consequence, and duration used in that decision.

The compiler rebuilds structural plans only when topology changes. Values, health evidence, session identity, validity, and bounds are refreshed at every step without treating routine runtime churn as a new neural architecture. This separation is important in the 15-second stress case, where value and session updates occur far more frequently than asset installation or removal.

\subsection{Constraint-Aware Local Bundle Construction}
\label{sec:timarl-local-feasibility}

Local hard constraints should not depend exclusively on a negative reward. An actor can learn to reduce violations, but safety-critical electrical and service limits require explicit semantics. \gls{timarl} keeps the learned preference in the raw bundle and applies a constraint-aware parameterisation to obtain the final local bundle.

Let $a^{\mathrm{raw}}_{i,t}$ denote the decoded bundle and $\mathcal{A}^{\mathrm{loc}}_{i,t}$ the feasible set induced by the local typed graph. The construction is

\begin{equation}
\bar{a}_{i,t}
=
\Phi_i
\left(
 a^{\mathrm{raw}}_{i,t},
 \mathcal{A}^{\mathrm{loc}}_{i,t}
\right),
\qquad
\bar{a}_{i,t}\in\mathcal{A}^{\mathrm{loc}}_{i,t}.
\label{eq:local-bundle-construction}
\end{equation}

Here $\Phi_i$ constructs the final bundle $\bar a_{i,t}$ from the raw decisions and the current local feasible set. The implementation uses an analytic joint projector. It first rejects a removed or invalid group, applies a forced typed fallback when one exists, clips the learned fraction to $[0,1]$, and respects the minimum operating power declared by the selected port. It then reserves mandatory deferrable starts and the \gls{ev} charging power needed to remain within the configured service schedule. Remaining charge and discharge requests are jointly scaled against total and phase-specific import or export headroom. A command that falls below its typed operating floor after scaling returns to \textsc{idle}.

For the running example, one local constraint is

\begin{equation}
P^{\mathrm{ch1}}_t
+
P^{\mathrm{ch2}}_t
+
P^{\mathrm{def}}_t
-
P^{\mathrm{bess,dis}}_t
\leq
P^{\max}_{A,t}.
\label{eq:running-local-constraint}
\end{equation}

Here $P^{\mathrm{ch1}}_t$ and $P^{\mathrm{ch2}}_t$ are the two charging powers, $P^{\mathrm{def}}_t$ is the deferrable load, and $P^{\mathrm{bess,dis}}_t$ is stationary-battery discharge. The quantity $P^{\max}_{A,t}$ is the remaining import headroom after accounting for the other loads and generation of building $A$. This total limit is evaluated together with phase incidence. If both chargers lie on L1, unused headroom on L3 cannot make their combination feasible. A three-phase service contributes to all three phase rows according to its declared incidence. The projector can retain a configurable headroom reserve and a completion margin for deferrables. This construction resolves competition within the member; \gls{rec} coordination remains with the learned policies.

Two diagnostics distinguish a useful feasibility mechanism from one that hides poor policy learning:

\begin{align}
I^{\mathrm{loc}}_t
&=\mathbb{1}\left[a^{\mathrm{raw}}_{i,t}\notin\mathcal{A}^{\mathrm{loc}}_{i,t}\right],\\
D^{\mathrm{loc}}_t
&=\left\|\bar{a}_{i,t}-a^{\mathrm{raw}}_{i,t}\right\|.
\end{align}

Here $I^{\mathrm{loc}}_t$ uses the indicator $\mathbb 1$ to record an infeasible raw bundle, and $D^{\mathrm{loc}}_t$ measures its adjustment. The norm is evaluated on the common physical-command vector after decoding; discrete mode changes are counted separately. Frequent or large adjustments indicate reliance on the projector and weak learning of coherent bundles. Requested, projected, and applied actions remain separate in the transition record, so a physically safe trajectory cannot conceal an infeasible raw policy. The optional shield-aware update removes an intervened group from the \gls{ppo} policy loss and can distil the final feasible decision into that group; unrelated groups remain eligible for learning.

\subsection{Shared Limits and Operational Supervision}
\label{sec:timarl-tcg}
\label{sec:timarl-shared-monitor}
\label{sec:timarl-runtime-assurance}

The \gls{tcg} records ownership, connection, dependency, and constraint relations. A charger contributes power to its building and to the phases on which it operates. These relations supply local construction and identify the shared resources affected by the resulting commands. The central set critic uses member summaries; the \gls{tcg} is the physical description supporting the decision process.

For buildings sharing a transformer, the \gls{rec} limit can be written as
\begin{equation}
\sum_{i\in\mathcal N_t}K_i^{\mathrm{import}}(\bar a_{i,t})
\leq P^{\max}_{\mathrm{transformer},t}.
\label{eq:running-shared-constraint}
\end{equation}
Here $K_i^{\mathrm{import}}$ maps member $i$'s constructed action to its resulting import under the current physical state, and $P^{\max}_{\mathrm{transformer},t}$ is the permitted aggregate import. The sum spans the active members $\mathcal N_t$. Satisfaction of each building's local limit alone does not establish this shared inequality. In the current method, \gls{rec} observations, coupled rewards, and centralised value estimation support learned coordination; residual shared violations remain part of the evaluation.

Operational supervision sits outside this learned coordination. The Fail-Safe Supervisor described in Chapter~\ref{chap:energAIze} validates equipment requests, limits them when necessary, and supplies a rapid rule-based fallback when inference is unavailable. Within \gls{timarl}, dependency closure can already disable an affected group or isolate a member whose essential measurements are missing. The execution record distinguishes those interface decisions, local projection, supervisory intervention, and the final equipment response. Later retraining can therefore retain the circumstances of an intervention without attributing the fallback's outcome to an action the actor never executed.

This division preserves a precise scope of authority. The local construction enforces the encoded member constraints when the feasible set is non-empty, the physical model is correct, and the procedure terminates normally. An infeasible service obligation is recorded as unmet; electrical protection cannot create missing capacity or charging time. \gls{rec}-wide hard enforcement would require a declared allocation or jointly certified constraint mechanism in addition to the local projector.

\section{Compatibility and Deployment}
\label{sec:timarl-compatibility-explainability}

\subsection{Interface Snapshots and Admission}
\label{sec:timarl-compatibility}

Every decision is associated with an immutable interface snapshot containing active entities, observation parts, action groups, ports, validity, bounds, dependencies, health, and schema versions. The snapshot allows an experiment transition to be reconstructed even when the runtime topology subsequently changes.

Before a checkpoint is restored, its compatibility signature is checked against the current contract, schema, type registry, health rules, learning backbone, critic, actor context, and credit-assignment mode. A new instance of a known charger type can be admitted without changing network architecture. A semantically unknown type is rejected. Compiler migration requires explicit opt-in and matching contract hashes; tensor compatibility alone is insufficient.

For a second installation, these checks identify whether the trained controller can be loaded as it stands or requires a revised interface and evaluation. It prevents a model from being resumed or moved merely because tensor dimensions happen to match. The packaging and deployment process in Chapter~\ref{chap:energAIze} uses this distinction at the model handoff.

\subsection{Checkpoint, Export, and Actor-Only Handoff}
\label{sec:timarl-persistence}

The \texttt{ti\_marl\_checkpoint\_v1} format stores the actor, critic, optional typed-group critic, optimisers, rollout, random-number states, compiler state, \gls{bc} state, resolved configuration, software versions, and compatibility signature. Checkpoint replacement is atomic. Restoring to a different compatible member composition is allowed because stable semantic types, rather than the active population, determine network shape.

Two exports have different purposes. The research artefact retains actor and critic state for evaluation and is marked \texttt{deployable: false}. A separate \texttt{ti\_marl\_deployment\_bundle\_v1} omits the critic and optimisation state and carries the actor, resolved typed interfaces, compiler and health rules, projector configuration, normalisation declaration, versions, and compatibility signature. This actor-only bundle defines the intended handoff to the inference lifecycle, but the current generic serving path does not yet certify it as a production \gls{timarl} deployment. Section~\ref{sec:energaize-contract} describes the general model contract and its relationship to method-specific runtime requirements.

\subsection{Execution-Level Decision Traces}
\label{sec:timarl-explainability}

Typed interfaces support explanations grounded in the actual decision path. A trace can identify

\begin{itemize}
  \item the active entities, modules, and schema versions;
  \item the local and \gls{rec} observations available to the actor;
  \item observations that were degraded, stale, missing, or failed;
  \item the action groups and ports that existed at that step;
  \item invalid ports and the dependency or health rule that invalidated them;
  \item the port and parameter selected for each active group;
  \item any local feasibility adjustment and the active local constraint;
  \item simulator execution feedback, including requested and applied values; and
  \item the final action executed by each asset.
\end{itemize}

For example, a trace can state that \texttt{charger\_1} selected \textsc{charge-ev} while \textsc{discharge-ev} was unavailable because the bound \gls{ev} \gls{soc} source was stale. It can separately record that \texttt{charger\_2} was scaled to preserve the L1 headroom and whether the simulator applied the resulting request. These statements follow compiler and constraint semantics; they are not inferred from attention weights.

Transitions are written to buffered gzip-compressed \gls{jsonl} chunks. Each transition references the current and next content-addressed snapshots. Complete snapshots are deduplicated and emitted periodically or when topology and health events require them. This keeps the audit trail useful on long high-frequency runs without serialising the complete interface graph at every step.

\subsection{Formal and Operational Properties}
\label{sec:timarl-properties}

The method supports several properties whose validity can be checked independently of control performance.

\paragraph{Entity-binding integrity.}
Every observation and action instance identifies its source, owner, and target entity. The compiler rejects ambiguous or orphaned bindings.

\paragraph{Interface soundness.}
A compiled valid port has an active owner module, satisfied hard dependencies, and a recognised parameter schema. Health-dependent contraction is deterministic under a fixed snapshot.

\paragraph{Local bundle feasibility.}
The final local bundle belongs to the encoded local feasible set when that set is non-empty, the construction terminates normally, and the encoded constraints describe the physical system correctly.

\paragraph{Information-boundary preservation.}
The deployed actor consumes only local interface objects and authorised aggregate observations. Privileged critic information is excluded from the actor package.

\paragraph{Safe admission.}
Unknown or incompatible schema elements do not become executable action ports merely because a numerical embedding can be produced.

\paragraph{Trace consistency.}
The explanation record references the same interface snapshot, selected ports, constraint rules, and final action used by execution.

These properties define testable invariants for compilation, admission, execution, and persistence. Energy and service performance are assessed through the experimental comparisons that follow.

\section{Experimental Methodology}
\label{sec:timarl-methodology}

The evaluation assesses the complete \gls{timarl} controller and its individual mechanisms through the four dimensions defined in Section~\ref{sec:intro-lexicon}. A controller must first provide useful energy control in the original installation. Resource changes and disturbances then test whether that performance can be retained, how much intervention is needed, and what work is required to continue operation. Chapters~\ref{chap:simulator} and~\ref{chap:control-benchmark} supply the common scenarios, annual comparison window, and scorecard.

\subsection{Reference Methods and Attribution}
\label{sec:timarl-comparison-design}

\gls{bau} with storage represents conventional operation, while \gls{rbc} Smart Local supplies a causal, inspectable control reference. The selected residual \gls{matd3} configuration supplies the practical learned reference developed in Chapter~\ref{chap:control-benchmark}. Transformer \gls{ppo} and \gls{matd3} examine reuse across known local resource instances. The \gls{milp} reference has a separate role: perfect-information scheduling indicates the opportunity within its declared optimisation model and horizon. Only a solver output replayed over the same scored window can enter the annual physical comparison.

The algorithmic comparison follows Section~\ref{sec:control-sota}. MOHITO-based and \gls{coin}-inspired adaptations examine coordination and modular reuse. Safety-layer and shielding references examine constraint enforcement, while \gls{marta}-based fault training examines robustness learned through adverse experience. \gls{gnn}-DT and HetGPS retain the relevant energy-specific comparisons. Each adaptation keeps its original mechanism explicit and declares the additional modelling, protection, and deployment support used in the common task.

The complete-controller comparison measures the annual operating outcome together with computation, integration effort, interventions, and recovery. Matched ablations then examine particular mechanisms. A protected entity-based \gls{mappo} reference holds the learner, actor observations, critic information, action scope, and training budget constant while changing representation and compilation. Further contrasts change protection, fault handling, or compatibility checks with the remaining configuration held fixed. These studies attribute the observed effects without narrowing the overall research problem to one mechanism.

\begin{table}[htbp]
\centering
\small
\caption{Matched comparisons for attributing \gls{timarl} behaviour.}
\label{tab:timarl-structural-baselines}
\begin{tabularx}{\textwidth}{@{}p{0.25\textwidth}p{0.36\textwidth}X@{}}
\toprule
\textbf{Comparison} & \textbf{Held constant} & \textbf{Question}\\
\midrule
Entity \gls{mappo} / \gls{timappo} & \gls{ppo}, information, population-aware critic, action scope, protection, budget. & What does the compiled interface add to a capable set representation?\\
\addlinespace
\gls{tippo} / \gls{timappo} & Compiler, actor, local construction, observations, reward. & What is gained through centralised value estimation?\\
\addlinespace
Health metadata / dependency closure & Detected events, measurement quality, thresholds, fallback. & Does changing dependent capabilities preserve more useful control?\\
\addlinespace
Protected / unprotected actor & Actor and held-out simulation trajectory. & How much feasibility is supplied by local construction?\\
\addlinespace
Local / action-conditioned decoding & Type vocabulary, observations, backbone, budget. & Does explicit group-to-observation binding improve repeated-asset decisions?\\
\addlinespace
Pure \gls{ppo} / prior knowledge & Evaluation inputs, action scope, later learning budget. & Which gains depend on demonstrations or causal auxiliaries?\\
\bottomrule
\end{tabularx}
\end{table}

Ablations of entity binding, semantic metadata, dynamic bounds, group credit, \gls{rec} observations, and compatibility checks extend this design. Each changes one declared mechanism. The simulator retains its physical protections when an algorithmic protection is removed, and both proposed and applied actions are recorded. This exposes the cost of an ablation without treating a rejected command as successful control.

\subsection{Annual Scenarios and Event Comparisons}
\label{sec:timarl-scenarios}

The principal suite contains Core-30 Normal, Topology, Safety, Robustness, and Combined, together with Premium-100 Normal and All-in. These are the seven annual scenarios used in Chapter~\ref{chap:control-benchmark}. Normal establishes the original installation; the focused variants separate resource changes, electrical pressure, and disturbances; the combined scenarios examine their interaction. Premium-100 contains 80 residential, 17 small-service, and 3 larger-service members. \gls{dr} is evaluated only where the scenario enables it, with its financial credit and physical additionality reported separately.

The first event comparison adds a known charger to an existing building. The assessment follows its admission, the binding of the new session observations, the additional action group, and the allocation of local headroom. Actor parameter count, compiler cost, required configuration work, and subsequent service are all recorded. A controller can accept the new group structurally yet schedule it poorly, so interface acceptance and control quality remain separate measurements.

The second comparison removes a measurement needed for one decision. A stale \gls{soc} tests whether dependent discharge is suspended while unaffected resources retain service. A failed main meter tests a broader loss of trustworthy headroom; loss of an aggregate \gls{rec} signal tests the continuation of local control. Candidate and reference receive the same actor-visible fault evidence. Recovery is timed against common freshness and hysteresis rules, distinguishing detection delay from the controller's response after detection.

The third comparison introduces another building. It follows admission of the local actor, population handling in the learner, continuity of the remaining members' trajectories, and the change in \gls{rec} outcomes. New members and new local assets are varied separately before being combined. Held-out compositions contain recognised types; an unknown type exercises rejection and fallback, with that event recorded as an admission decision.

Each annual scorecard includes the complete controller--supervisor trajectory. Event-centred records explain what happened within it: learned decisions retained, groups suspended, fallback duration, recovery, and missed service. Preparation or retraining for another installation is measured outside the frozen test year. This keeps the full operating consequence visible while preventing a change of policy during evaluation from being presented as immediate reuse.

\subsection{Measurements across the Four Research Dimensions}
\label{sec:timarl-measures}

The absolute scorecard retains energy-purchase cost, emissions, \gls{rec} import and export, \gls{pv} self-consumption, daily and maximum peak, upward ramping, \gls{ev} minimum service and target accuracy, electrical-violation energy, action correction, battery throughput, \gls{v2g} energy, and the number of members benefiting. Definitions, units, denominators, and aggregation follow Chapter~\ref{chap:simulator}. A single reward value cannot express all these outcomes, particularly when a low cost accompanies missed departures or frequent intervention.

\begin{table}[htbp]
\centering
\small
\caption{Mechanisms and complementary evidence for the four research dimensions.}
\label{tab:timarl-four-dimensions}
\begin{tabularx}{\textwidth}{@{}p{0.18\textwidth}p{0.37\textwidth}X@{}}
\toprule
\textbf{Dimension} & \textbf{Mechanisms examined} & \textbf{Evidence retained}\\
\midrule
Scalability and dynamic reconfiguration & Shared actor, variable observation/group sets, central set critic, structural caching. & Energy and service across counts; training/adaptation time, memory, throughput, compiler and inference latency.\\
\addlinespace
Deployability & Versioned interfaces, compatibility checks, actor bundle, stable resource bindings. & Configuration and validation effort; accepted/rejected changes; replay agreement, loading and recovery time.\\
\addlinespace
Safety & Dependency validity, physical bounds, local construction, external supervision. & Requested/projected/applied commands; violated energy, correction frequency and magnitude, missed service.\\
\addlinespace
Robustness & Scoped capability loss, surviving local decisions, fallback, recovery rules. & Fault exposure, unavailable groups, retained decisions, failed events, fallback and recovery duration, service deficits.\\
\bottomrule
\end{tabularx}
\end{table}

The same mechanism can affect several dimensions. Reusing an actor after charger installation can support dynamic reconfiguration while reducing training, adaptation, and integration effort. Computational growth is assessed under scalability; the effort needed to prepare and maintain that actor belongs to deployability. Health-dependent port removal contributes to safety by excluding an unsupported request and to robustness by retaining unaffected decisions. The measurements in Table~\ref{tab:timarl-four-dimensions} give those related claims distinct empirical meanings.

Computational results identify hardware, batch size, member count, observation count, active groups, time resolution, and deterministic decoding settings. Constant parameter count is reported alongside the increasing work of processing additional resources. Operational measurements similarly identify their scope: validating an actor bundle in software and commissioning a real-world installation involve different tasks and durations.

\subsection{Training, Selection, and Paired Evaluation}
\label{sec:timarl-protocol}

The temporal split follows the annual protocol of Chapter~\ref{chap:control-benchmark}. Training uses 2023 excluding the first fourteen \gls{utc} days of February, May, August, and November, which are reserved for validation. Episode boundaries prevent a learning target from crossing into a held-out block. Preprocessing statistics are fitted on the training inputs and frozen. The complete 2024 sequence supplies the common test window, with no policy updates or tuning against its outcomes.

The repetition design uses development seeds 123, 456, and 789 and five independently trained confirmation replicas with seeds 1009, 2027, 3037, 4051, and 5059. Candidate and reference share the same scenario chronology, initial physical conditions, exogenous inputs, and event schedule. Neural seeds vary the learned controller rather than the physical task. The principal estimate is the paired annual difference for each seed, with individual seed outcomes and the uncertainty interval retained. Daily or seasonal breakdowns interpret that complete window without replacing it.

Nominal selection first checks electrical and service requirements. Applied electrical-violation energy has a numerical tolerance of $10^{-6}$~kWh; at least 99\% of feasible \gls{ev} departures must satisfy the minimum service requirement; appliance deadlines and unserved energy are checked separately. \gls{ev} target accuracy and appliance service must not deteriorate relative to the matched reference. Daily peak and ramping may increase by at most 5\%, while \gls{pv} self-consumption may decrease by at most two percentage points. These thresholds are declared selection rules, applied before the test year is opened.

Among candidates passing those conditions, the selected policy must reduce cost by at least the larger of EUR~0.01 and 0.1\% of the absolute reference cost. Daily peak, ramping, and solar utilisation resolve ties. In disturbed or physically infeasible periods, the record instead retains exposure, requested service, delivered service, and recovery. A power outage without backup cannot be judged by an uninterrupted-supply requirement. The appropriate comparison uses the same event for both controllers and a matched undisturbed replay to explain the attributable loss.

\section{Results and Comparative Analysis}
\label{sec:timarl-results}

The available results establish the operation of the implemented interface, learner, and persistence mechanisms. The annual experimental comparisons specified above remain to be completed. This section reports the software evidence directly and retains the quantitative control questions separately, without assigning performance values to unexecuted comparisons.

\subsection{Interface and Learning Implementation}
\label{sec:timarl-results-nominal}

The focused regression suite passed all 139 tests in the repository snapshot based on Algorithms commit \texttt{a3c8194e}, using \texttt{softcpsrecsimulator==1.8.0}. The retained test log and source hashes identify the tested state. The suite covers compilation, policy runtime, checkpoints, traces, the public typed-interface path, a dynamic simulator integration case, and a high-frequency topology fixture. Table~\ref{tab:timarl-verified-evidence} summarises the exercised behaviour.

\begin{table}[htbp]
\centering
\small
\caption{\gls{timarl} behaviour established by the focused 139-test regression suite.}
\label{tab:timarl-verified-evidence}
\begin{tabularx}{\textwidth}{@{}p{0.25\textwidth}p{0.39\textwidth}X@{}}
\toprule
\textbf{Mechanism} & \textbf{Exercised behaviour} & \textbf{Interpretation}\\
\midrule
Interface compilation & Determinism, ordering, semantic distinction, unknown-type rejection, dynamic bounds, repeated groups. & The tested resource associations survive changes in representation order and known instances.\\
\addlinespace
Local construction & Total/per-phase limits, \gls{ev} service floors, discharge reserves, operating minima, mandatory starts. & The tested bundles follow their encoded local contracts, with proposal and adjustment retained.\\
\addlinespace
Local and central training & \gls{tippo} and \gls{timappo} execute training in the simulator integration case. & Both critic arrangements operate with the common typed actor.\\
\addlinespace
Changing population & Join/leave trajectory boundaries, compatible restoration, unchanged parameter dimensions. & Population changes preserve the tested member identities and surviving trajectories.\\
\addlinespace
Health and execution records & Fault evidence reaches health, validity, closure records, commands, and execution feedback. & The tested intervention can be traced to its affected capability and rule.\\
\addlinespace
High-frequency topology & Six structural events yield seven structural versions; the compiler p95 meets a 500~ms test ceiling. & Routine value updates avoid unnecessary structural rebuilding in the 15-second fixture.\\
\bottomrule
\end{tabularx}
\end{table}

The integration tests establish a concrete connection between the formal objects and the simulator. A resource change reaches the compiler, creates or removes the corresponding groups, and passes through actor evaluation and action encoding. The topology fixture also distinguishes structural rebuilding from ordinary measurement updates. Its latency ceiling is a regression condition for that fixture; deployment timing requires the hardware-specific measurements in Section~\ref{sec:timarl-measures}.

\subsection{Comparative Questions and Evidence Scope}
\label{sec:timarl-results-reconfiguration}
\label{sec:timarl-results-health}
\label{sec:timarl-results-coordination}
\label{sec:timarl-results-scale}
\label{sec:timarl-results-ablations}

The tests support the implemented treatment of identity, validity, local construction, and restoration. The remaining empirical question is the benefit obtained when those mechanisms are used throughout an operating year. For the added charger, that benefit would appear in retained service and reduced rebuilding or retraining effort. For the missing measurement, it would appear in fewer unnecessary suspensions while respecting the affected physical constraints. For the joining building, it would appear in useful coordination at an acceptable computational and integration cost.

Those outcomes require the matched comparisons in Section~\ref{sec:timarl-methodology}. A protected entity-based controller may achieve comparable energy performance; an explicit contract can still reduce integration ambiguity, but that becomes a deployability result requiring its own measurement. Conversely, accepting another charger without rebuilding the network establishes structural reuse while leaving its learned schedule to be assessed. The illustrative tables in Chapter~\ref{chap:control-benchmark} supply a reporting layout, not evidence of an advantage for \gls{timarl}.

\section{Discussion}
\label{sec:timarl-discussion}

The scientific question is whether \gls{timarl} sustains useful \gls{rec} control across scalability, deployability, safety, and robustness. The method connects learned coordination to resource interpretation, local protection, fault response, and model handoff. Its comparison with the state of the art therefore follows both the complete operating result and the contribution of each mechanism.

\subsection{Relationship to State-of-the-Art Methods}
\label{sec:timarl-neighbour-comparison}

MOHITO's task-open formulation is especially relevant when service demands appear and disappear. In the running example, a vehicle arrival creates a session with an energy requirement and a deadline. \gls{timarl} distinguishes that session from the charger that serves it and from the building that owns the local controller. Its compiled action group remains attached to the charger while the target session changes. The distinction becomes consequential when a charger is installed or fails: that event changes a capability and its dependencies, rather than merely adding or removing a task. The energy adaptation in Section~\ref{sec:control-variable-references} makes these objects explicit before comparing the methods.

\gls{coin} provide a complementary form of modularity. Their separation of perception, a shared controller, and action modules concerns which learned components can be reused across interface families. \gls{timarl} places explicit runtime bindings and constraints around repeated instances of recognised types. Two identical charger types share learned processing while retaining separate sessions, bounds, and phase incidence. This identifies the additional relationship under study; a modular reference equipped with the same bindings and protection is a stronger comparator than an unprotected fixed vector.

The same reasoning applies to graph-based energy methods. \gls{gnn}-DT associates actions with charging resources across changing fleets, and HetGPS combines a reusable graph policy with physical correction. These are substantial precedents for the representation and safety components of the problem. \gls{timarl}'s distinguishing design is the common executable interface used by the actor, dependency rules, local construction, compatibility checks, and traces. Its empirical contribution depends on whether this connection improves the outcomes or effort measured under matched conditions.

Fault-tolerant learning supplies a further algorithmic comparison. \gls{marta} changes the adverse experience used in training, whereas \gls{timarl}'s dependency rules determine which capabilities remain available after the recorded evidence changes. The two mechanisms can complement one another. A matched experiment compares ordinary and fault-aware training under the same detection, protection, and recovery conditions, so that retained service can be attributed to the learner or the operational response.

\begin{table}[htbp]
\centering
\small
\caption{Relationships to neighbouring mechanisms and the specific \gls{timarl} comparison.}
\label{tab:timarl-progression}
\begin{tabularx}{\textwidth}{@{}p{0.22\textwidth}p{0.36\textwidth}X@{}}
\toprule
\textbf{Neighbouring mechanism} & \textbf{Capability already established} & \textbf{\gls{timarl} comparison}\\
\midrule
Task-open hypergraphs & Relations between changing tasks, agents, and actions. & Persistent equipment and temporary sessions, with capability-dependent validity and bounds.\\
\addlinespace
\gls{coin} & Recombination of perception and action modules around a shared controller. & Runtime bindings of repeated resource instances and their constraints through one contract.\\
\addlinespace
Energy graph policies & Resource-level decisions, fleet/topology transfer, and, in HetGPS, corrective control. & Equal physical information and protection, with compilation and operational effort measured separately.\\
\addlinespace
Safety layers and compositional shields & Explicit correction or certified local obligations supporting declared safety properties. & Local feasible construction linked to current resources; global guarantees require shared assumptions and enforcement.\\
\addlinespace
Fault-tolerant learning & Selective adverse experience improves the policy response to the studied faults. & Equal fault exposure, detection and fallback; retained service, recovery, and training cost.\\
\addlinespace
Model lifecycle tooling & Versioning, packaging, validation, and deployment management. & Compatibility of the complete resource-to-command interpretation at the model handoff.\\
\bottomrule
\end{tabularx}
\end{table}

\subsection{Scalability and Deployability}

A shared decoder allows the third charger to use the same learned parameters as the first two. The set critic similarly accepts another compatible building without extending a concatenated input layer. These choices address structural growth. Their practical value depends on processing time, memory, retained service, and the training or adaptation needed after the change. More resources still produce more observations, groups, and interactions to evaluate.

Deployability asks what must be done to use that controller in the changed installation. The new charger must be identified, associated with the right measurements, assigned to its electrical constraints, and checked against the model's supported types. The compiled interface and compatibility signature preserve these associations in the actor bundle. They make reuse explicit enough to validate, and provide a common basis for the services described in Chapter~\ref{chap:energAIze}.

The two dimensions overlap without becoming interchangeable. Avoiding network reconstruction can reduce computational work and installation effort at the same time. Yet a small model may still require extensive manual mapping, and a well-packaged model may be too expensive to evaluate on the target hardware. Separate timing, resource, and integration measurements reveal which benefit the method provides. The current actor bundle establishes the algorithm-side handoff; complete production serving and commissioning remain additional evidence requirements.

\subsection{Safety and Robustness}

Safety concerns the action that reaches the resource and its consequences. A stale vehicle \gls{soc} can remove the discharge port because its reserve cannot be established. A valid charging request can still be reduced when two chargers compete for the same phase headroom. These are distinct interventions: dependency closure determines available capabilities, while local construction reconciles their combined physical effect. Requested, projected, and applied records make both visible alongside missed departures and other service outcomes.

Formal shielding provides an important comparison. Joint safety can follow from local obligations when their conjunction has been certified against the global specification \parencite{adalat2026contracts}. \gls{timarl}'s member-level projector has a narrower authority: it enforces the encoded local set. Shared transformer safety remains an outcome of the learned joint operation and the declared external protections. The four-dimension assessment therefore cannot convert local feasibility into a claim of complete \gls{rec} safety.

Robustness concerns how much useful control remains during the same disturbance and how operation recovers. Disabling one charger's discharge while retaining another charger's service is more selective than suspending the entire building. That selectivity depends on trustworthy fault evidence and correct dependencies. Training against corrupted inputs can complement it by improving decisions over the remaining information. Equal detection thresholds and fallback rules are essential to determine whether an observed improvement comes from the policy, the compiled response, or earlier knowledge of the fault.

\subsection{Scope of the Contribution}

\gls{timarl} integrates reusable multi-agent learning with local protection, health-dependent operation, and a compatible model handoff. The common typed description gives these mechanisms a consistent account of resources, observations, available actions, and constraints. Their combined value is assessed through retained energy and service performance, computation, integration effort, intervention, and recovery. Matched comparisons determine which mechanisms account for that value.

The present implementation supports recognised resource and action families, declared health rules, and encoded local constraints. Its mean-pooled critic summarises the active \gls{rec}, while its actor bundle defines the algorithm-side deployment requirements. The focused tests establish operation under the exercised conditions. Comparative control and deployment evidence determines the practical benefit, including cases where another method achieves similar performance with different preparation or computational costs.

\section{Chapter Summary}
\label{sec:timarl-summary}

\gls{timarl} connects multi-agent learning, local protection, health-dependent operation, and model handoff through a common typed description. Shared actors and local or central critics support \gls{rec} control, while resource bindings and constraints preserve the meaning of decisions. The assessment follows scalability, deployability, safety, and robustness through the same energy and service task, with matched comparisons examining the contribution of individual mechanisms.

Chapter~\ref{chap:real-world} returns to the installations and operational experience behind these requirements. Within the architecture of Chapter~\ref{chap:energAIze}, data services, learned decisions, supervision, and equipment interfaces must sustain that control together.

\postdisplaypenalty=\timarlpostdisplaypenalty\relax
